\section{Introduction}\label{sec:intro}

A trading desk may solve the same multi-asset PDE many times while repricing a book. The wall-clock cost of one solve therefore affects the throughput available for valuation and risk analysis. Niesen and Wright's option-pricing working paper reported an approximately 700\,s historical CPU result for one three-asset configuration.\footnote{J.~Niesen and W.~M.~Wright, \emph{A Krylov Subspace Method for Option Pricing}, working paper, DOI: \href{https://doi.org/10.2139/ssrn.1799124}{10.2139/ssrn.1799124}.} This figure provides historical context, but it is not a contemporary performance baseline.

In the implementation studied here, a Krylov exponential integrator spends over 70\% of its time in sparse matrix--vector products (SpMV), which move far more bytes than they compute. Most of the remaining time is spent in an orthogonalization whose inner products require reductions. The first cost is data movement through the memory hierarchy; the second is synchronization between workers \cite{ballard-2011,hoemmen-2010}. Communication-avoiding (CA) methods target both costs. A matrix-powers kernel reuses an operator tile across several Krylov vectors \cite{demmel-2008,mohiyuddin-2009}, while grouped orthogonalization reduces the number of synchronization phases \cite{hoemmen-2010,carson-2015}. These savings come with extra arithmetic and storage.

Whether this exchange pays depends on the operator, grid, implementation, and machine. This thesis first identifies where communication avoidance has an opportunity to help and then tests whether the implementation realizes that opportunity on the available hardware.

Three artifacts support that study. A \textbf{benchmark pricer} in C++23 implements five methods (Crank--Nicolson, two ADI variants, a one-shot matrix exponential, and a Krylov exponential integrator) on two three-asset European options and isolates temporal error for a fixed semi-discretization (\Cref{sec:benchmark}). A \textbf{regime study} (\Cref{sec:regime}) builds a dimensionless map of the CA design space. Its axes normalize demand by machine-dependent supply; their unit lines are opportunity indicators rather than promises of speedup or hardware portability. A \textbf{distributed CA exponential integrator} (\Cref{sec:integrator}) then evaluates those mechanisms on up to four GPUs across two nodes and asks whether their savings survive contact with a complete adaptive solver.

One design rule keeps the map non-circular: its coordinates are computed before a run from the problem size, sparsity pattern, spectrum, and independently measured machine constants. Wall-clock time, orthogonality loss, and achieved traffic do not enter the calculation. The experiment can therefore test the coordinates as opportunity indicators and leave the measured crossover as a result.

The thesis proceeds as follows. \Cref{sec:model} states the pricing problem and its discretization. \Cref{sec:ksm} presents the Krylov exponential integrator, its cost structure, and the CA reformulations that target it. \Cref{sec:platform} describes the experimental platform and measurement methodology. \Cref{sec:benchmark} follows the historical Niesen--Wright option-pricing experiment as a temporal-integration control. \Cref{sec:profile} profiles the solver and places its kernels on a cache-aware roofline. \Cref{sec:scaling} presents the OpenMP scaling study that motivates the two CA mechanisms. \Cref{sec:regime} develops the regime map and its numerical and measured results. \Cref{sec:gpu} recomputes the coordinates for a GPU memory system and tests their transfer across two distinct GPU architectures. \Cref{sec:integrator} builds the distributed integrator and measures topology-specific performance on a full solve. \Cref{sec:future} names what the measurements leave open, and \Cref{sec:conclusion} concludes.
